% math.tex — geom module
% Mathematical results supporting crates/src/geom/
%
% Labels defined here:
%   def:symplectic-form, def:J0, def:ehz-capacity, def:lagrangian-product,
%   def:symplectic-product, def:cross-product-4d, def:face-lattice,
%   def:polar-body, def:polytope-dual, def:polygon-h-rep, def:polygon-area,
%   def:reeb-vector-field, def:volume, def:systolic-ratio,
%   lem:piecewise-linear-reeb, lem:positive-span, lem:rational-pipeline,
%   lem:vertex-enumeration, lem:shoelace,
%   thm:hko-counterexample,
%   prop:capacity-symplectic-product
%
% Sources: thesis/basic-definitions.tex (verified content),
%   .rs doc comments (stubs for labels not yet in thesis)

\documentclass{article}
\usepackage{amsmath,amssymb,amsthm}
\newtheorem{theorem}{Theorem}
\newtheorem{lemma}[theorem]{Lemma}
\newtheorem{definition}[theorem]{Definition}
\newtheorem{corollary}[theorem]{Corollary}
\newtheorem{proposition}[theorem]{Proposition}
\newtheorem{remark}[theorem]{Remark}
\begin{document}

% ====================================================================
% Symplectic structure
% ====================================================================

\begin{definition}[Standard symplectic structure]\label{def:symplectic-form}
% Jörn: text approved (1aac085)
Write $\mathbb{R}^4$ with coordinates $(q_1, q_2, p_1, p_2)$.
The \emph{standard complex structure} is
\[
  J_0 = \begin{pmatrix} 0 & -I_2 \\ I_2 & 0 \end{pmatrix},
  \qquad\text{i.e.,}\quad
  J_0(q, p) = (-p, q),
\]
where $I_2$ is the $2 \times 2$ identity matrix.
The \emph{standard symplectic form} is
\[
  \omega_0(u, v) \coloneqq \langle J_0 u, v \rangle
  = \sum_{i=1}^{2} (u_{q_i} v_{p_i} - u_{p_i} v_{q_i}),
\]
equivalently $\omega_0 = dq_1 \wedge dp_1 + dq_2 \wedge dp_2$.
\end{definition}


\begin{definition}[Standard complex structure $J_0$]\label{def:J0}
% [TODO: JÖRN - verify: J0 is defined within def:symplectic-form above.
%   This label exists as a separate cross-reference target for code that
%   references J0 specifically (e.g. symplectic_form.rs). Consider whether
%   this should remain a separate label or be merged into def:symplectic-form.]
The standard complex structure $J_0 \in \mathbb{R}^{4 \times 4}$ is
\[
  J_0 = \begin{pmatrix} 0 & -I_2 \\ I_2 & 0 \end{pmatrix},
\]
satisfying $J_0^2 = -I_4$, $J_0^{-1} = J_0^\top = -J_0$,
$\omega_0(J_0 u, J_0 v) = \omega_0(u,v)$, and $|J_0 u| = |u|$.
In the 2D case, $J_0^{(2)} = \bigl(\begin{smallmatrix} 0 & -1 \\ 1 & 0 \end{smallmatrix}\bigr)$.
\end{definition}


% ====================================================================
% Capacity and systolic ratio
% ====================================================================

\begin{definition}[EHZ capacity]\label{def:ehz-capacity}
% Jörn: text approved (1aac085)
The \emph{Ekeland--Hofer--Zehnder capacity} of a convex body $K$ with smooth
boundary is
\[
  c_{\mathrm{EHZ}}(K) \coloneqq \min\bigl\{ A(\gamma) :
  \gamma \text{ is a closed characteristic on } \partial K \bigr\}.
\]
The minimum exists and is positive.
The definition extends to all convex bodies by approximation;
for polytopes, the minimum is attained by a generalized Reeb orbit.
\end{definition}


\begin{definition}[Systolic ratio]\label{def:systolic-ratio}
% [TODO: JÖRN - verify statement]
% Source: dataset.rs doc comment
The \emph{systolic ratio} of a convex body $K \subset \mathbb{R}^4$ is
\[
  \operatorname{sys}(K) \coloneqq \frac{c_{\mathrm{EHZ}}(K)^2}{2\,\operatorname{vol}(K)}.
\]
Viterbo's conjecture asserts $\operatorname{sys}(K) \leq 1$ for all convex bodies $K$.
\end{definition}


% ====================================================================
% Products
% ====================================================================

\begin{definition}[Lagrangian product]\label{def:lagrangian-product}
% Source: thesis/lagrangian-product-algorithm-proof.tex
Let $K_q \subset L_q$ and $K_p \subset L_p$ be convex polygons
with $0 \in \operatorname{int}(K_q)$ and $0 \in \operatorname{int}(K_p)$,
where $L_q = \{(q_1, q_2, 0, 0)\}$ and $L_p = \{(0, 0, p_1, p_2)\}$
are complementary Lagrangian subspaces.
The \emph{Lagrangian product} is the polytope
\[
  K \coloneqq K_q \times K_p
  = \{ (q, p) \in \mathbb{R}^4 : q \in K_q,\; p \in K_p \}.
\]
\end{definition}


\begin{definition}[Symplectic product]\label{def:symplectic-product}
% [TODO: JÖRN - verify statement]
% Source: known_polytopes.rs doc comment
Let $A, B \subset \mathbb{R}^2$ be convex bodies containing
the origin in their interiors, where we view the first copy
of~$\mathbb{R}^2$ as the symplectic plane $(q_1, p_1)$ and
the second as $(q_2, p_2)$.
The \emph{symplectic product} is
\[
  A \times_S B
    \coloneqq \{(q_1, q_2, p_1, p_2) \in \mathbb{R}^4
      : (q_1, p_1) \in A,\; (q_2, p_2) \in B\}.
\]
\end{definition}


\begin{proposition}[Capacity of symplectic products]%
\label{prop:capacity-symplectic-product}
% Source: thesis/basic-definitions.tex
Let $A, B \subset \mathbb{R}^2$ be convex bodies containing
the origin in their interiors, viewed as in
Definition~\ref{def:symplectic-product}.
Then
\[
  c_{\mathrm{EHZ}}(A \times_S B)
    = \min\!\bigl(\operatorname{area}(A),\;
                   \operatorname{area}(B)\bigr).
\]
\end{proposition}

\begin{proof}
We reduce to a product of disks via an area-preserving
map, then squeeze between a ball and a cylinder.

On each copy of~$\mathbb{R}^2$, the symplectic form
restricts to $dq_i \wedge dp_i$, which is the standard
area form, so area-preserving diffeomorphisms are
symplectomorphisms.

We first show $c_{\mathrm{EHZ}}(A) = \operatorname{area}(A)$
for every convex body~$A$ in~$\mathbb{R}^2$.
By the theorem of Dacorogna and
Moser, for any two bounded
domains~$\Omega_1, \Omega_2 \subset \mathbb{R}^2$ with
smooth boundaries and equal area, there exists an
area-preserving diffeomorphism
$\varphi \colon \overline{\Omega}_1 \to \overline{\Omega}_2$.
The construction solves a boundary-value problem whose
solution vanishes on~$\partial\Omega_1$, so the resulting
divergence-free vector field extends by zero outside
$\Omega_1$; its time-$1$ flow is a global area-preserving
(hence symplectic) diffeomorphism of~$\mathbb{R}^2$.
Taking $\Omega_1$ to be a smooth convex body~$A$
and $\Omega_2 = B^2(r)$ with $\pi r^2 = \operatorname{area}(A)$,
symplectic invariance gives
$c_{\mathrm{EHZ}}(A) = c_{\mathrm{EHZ}}(B^2(r)) = \pi r^2
= \operatorname{area}(A)$.
For non-smooth convex bodies, let
$A^{-\varepsilon} \subset A \subset A^{+\varepsilon}$
be smooth inner and outer parallel bodies;
monotonicity and
$\operatorname{area}(A^{\pm\varepsilon}) \to
\operatorname{area}(A)$ yield
$c_{\mathrm{EHZ}}(A) = \operatorname{area}(A)$.

Now let $r_A, r_B > 0$ satisfy
$\pi r_A^2 = \operatorname{area}(A)$ and
$\pi r_B^2 = \operatorname{area}(B)$.
The above gives global symplectomorphisms
$\varphi_1, \varphi_2 \colon \mathbb{R}^2 \to \mathbb{R}^2$
mapping $A$ to~$B^2(r_A)$ and $B$ to~$B^2(r_B)$
respectively.
Since $\omega_0 = dq_1 \wedge dp_1 + dq_2 \wedge dp_2$
decomposes over the two planes, the product
$\varphi_1 \times \varphi_2$ is a symplectomorphism
of~$\mathbb{R}^4$.
By symplectic invariance,
\[
  c_{\mathrm{EHZ}}(A \times_S B)
    = c_{\mathrm{EHZ}}\bigl(B^2(r_A) \times_S B^2(r_B)\bigr).
\]
Assume without loss of generality that $r_A \leq r_B$.
Then
\[
  B^4(r_A)
    \;\subset\; B^2(r_A) \times_S B^2(r_B)
    \;\subset\; B^2(r_A) \times \mathbb{R}^2,
\]
where the first inclusion holds because $r_A \leq r_B$
and the second because $B^2(r_B) \subset \mathbb{R}^2$.
By monotonicity and normalization,
\[
  \pi r_A^2
    = c_{\mathrm{EHZ}}(B^4(r_A))
    \leq c_{\mathrm{EHZ}}\bigl(B^2(r_A) \times_S B^2(r_B)\bigr)
    \leq c_{\mathrm{EHZ}}(B^2(r_A) \times \mathbb{R}^2)
    = \pi r_A^2.
\]
Hence $c_{\mathrm{EHZ}}(A \times_S B)
  = \pi r_A^2
  = \min(\operatorname{area}(A), \operatorname{area}(B))$.
\end{proof}


% ====================================================================
% Polytope geometry
% ====================================================================

\begin{definition}[Polytope dual / polar body]\label{def:polytope-dual}
% [TODO: JÖRN - verify statement]
% Source: polytope.rs doc comment
Let $K \subset \mathbb{R}^4$ be a polytope with $0 \in \operatorname{int}(K)$
given in $H$-representation $K = \{x : \langle a_i, x\rangle \leq 1\}_{i=1}^F$.
The \emph{polar body} (or \emph{dual}) is
\[
  K^\circ \coloneqq \{y \in \mathbb{R}^4 : \langle x, y \rangle \leq 1
    \text{ for all } x \in K\}
  = \operatorname{conv}(a_1, \ldots, a_F).
\]
\end{definition}

\begin{remark}\label{rem:polar-body-note}
% This remark connects def:polytope-dual to the separate label def:polar-body
% used in some .rs files. They refer to the same concept.
The label \texttt{def:polar-body} in some \texttt{.rs} files
refers to the same concept as Definition~\ref{def:polytope-dual}.
\end{remark}
\label{def:polar-body}% alias for cross-references


\begin{definition}[Face lattice / skeleton]\label{def:face-lattice}
% [TODO: JÖRN - verify statement]
% Source: skeleton.rs doc comment
The \emph{skeleton} (or \emph{face lattice}) of a 4D convex polytope $K$
is the combinatorial structure recording all faces of $K$
(vertices, edges, 2-faces, facets) and their incidence relations.
Two facets $F_i, F_j$ are \emph{adjacent} if $F_i \cap F_j$ is a
face of dimension $\geq 1$.
\end{definition}


\begin{definition}[Cross product in $\mathbb{R}^4$]\label{def:cross-product-4d}
% [TODO: JÖRN - verify statement]
% Source: cross_product_4d.rs doc comment
The \emph{cross product} of three vectors $u, v, w \in \mathbb{R}^4$
is the unique vector $u \times v \times w \in \mathbb{R}^4$
perpendicular to all three, with magnitude equal to the volume
of the parallelepiped spanned by $u, v, w$, and orientation
determined by the right-hand rule.
Equivalently,
\[
  (u \times v \times w)_i
  = (-1)^{i+1} \det\bigl(M_i\bigr),
\]
where $M_i$ is the $3 \times 3$ matrix obtained by deleting
row~$i$ from the $4 \times 3$ matrix $(u \mid v \mid w)$.
\end{definition}


% ====================================================================
% Polygon geometry
% ====================================================================

\begin{definition}[Convex polygon $H$-representation]\label{def:polygon-h-rep}
% [TODO: JÖRN - verify statement]
% Source: polygon.rs doc comment
A 2D \emph{convex polygon in $H$-representation} is specified by
outward unit normals $n_1, \ldots, n_m \in S^1$ and positive
heights $h_1, \ldots, h_m > 0$:
\[
  P = \bigcap_{i=1}^{m} \{x \in \mathbb{R}^2 : \langle n_i, x \rangle \leq h_i\}.
\]
\end{definition}


\begin{definition}[Polygon area]\label{def:polygon-area}
% [TODO: JÖRN - verify statement]
% Source: polygon.rs doc comment
The \emph{area} of a convex polygon $P$ given in $H$-representation
(Definition~\ref{def:polygon-h-rep}) is computed via triangulation
from the centroid or via the shoelace formula on the vertex list.
\end{definition}


% ====================================================================
% Volume
% ====================================================================

\begin{definition}[Volume of a 4D polytope]\label{def:volume}
% [TODO: JÖRN - verify statement]
% Source: volume.rs, qhull.rs doc comments
The \emph{volume} of a 4D convex polytope $K$ is the standard
Lebesgue measure $\operatorname{vol}(K) = \int_K dx$.
In the implementation, volume is computed by triangulating $K$
into simplices via qhull and summing the simplex volumes:
\[
  \operatorname{vol}(\Delta)
  = \frac{1}{4!}\, |\det(v_1 - v_0 \mid v_2 - v_0 \mid v_3 - v_0 \mid v_4 - v_0)|.
\]
\end{definition}


% ====================================================================
% Reeb dynamics on polytopes
% ====================================================================

\begin{definition}[Reeb vector field on polytopes]\label{def:reeb-vector-field}
% [TODO: JÖRN - verify statement]
% Source: reeb_trajectory.rs doc comment
On a facet $F_i$ of a polytope $K$ with outward unit normal $n_i$
and height $h_i$, the \emph{Reeb vector} is
\[
  R_i = \frac{2}{h_i}\, J_0\, n_i.
\]
The Reeb vector $R_i$ is tangent to $F_i$
($\langle n_i, R_i \rangle = 0$, since
$\langle n_i, J_0 n_i \rangle = \omega_0(n_i, n_i) = 0$).
\end{definition}


\begin{lemma}[Piecewise-linear Reeb trajectories]\label{lem:piecewise-linear-reeb}
% [TODO: JÖRN - verify statement, add proof]
% Source: reeb_trajectory.rs doc comment
Reeb trajectories on polytope boundaries are piecewise linear:
on each facet $F_i$, the trajectory moves with constant velocity $R_i$.
At a breakpoint where the trajectory transitions from $F_i$ to $F_j$,
the velocity changes discontinuously from $R_i$ to $R_j$.
\end{lemma}


% ====================================================================
% Computational geometry lemmas
% ====================================================================

\begin{lemma}[Positive span criterion]\label{lem:positive-span}
% [TODO: JÖRN - verify statement, add proof]
% Source: validation.rs, vertex_enumeration.rs doc comments
The outward normals $n_1, \ldots, n_F$ of a polytope $K \subset \mathbb{R}^4$
positively span $\mathbb{R}^4$ if and only if the recession cone
$\operatorname{rec}(K) = \{0\}$, i.e., $K$ is bounded.
This is used to verify that a given $H$-representation defines
a bounded polytope.
\end{lemma}


\begin{lemma}[Exact vertex enumeration]\label{lem:vertex-enumeration}
% [TODO: JÖRN - verify statement, add proof]
% Source: vertex_enumeration.rs doc comment
For a 4D polytope $K$ given by $F$ halfspaces, the vertices
are obtained by solving $\langle n_{i_k}, x \rangle = h_{i_k}$
for all $\binom{F}{4}$ four-element subsets $\{i_1, i_2, i_3, i_4\}$,
then filtering for points satisfying all remaining inequalities.
The implementation uses exact rational arithmetic to avoid
numerical errors in degeneracy detection.
\end{lemma}


\begin{lemma}[Rational arithmetic pipeline]\label{lem:rational-pipeline}
% [TODO: JÖRN - verify statement, add proof]
% Source: rational_arithmetic.rs doc comment
All discrete/combinatorial decisions in the geometry pipeline
(vertex adjacency, symplectic sign $\omega_0(n_i, n_j) \gtrless 0$,
halfspace irredundancy) are computed using exact rational arithmetic
to prevent numerical misclassification.
\end{lemma}


% ====================================================================
% Shoelace formula
% ====================================================================

\begin{lemma}[Symplectic shoelace formula]\label{lem:shoelace}
% Jörn: text approved (b7a8bc0) — from \begin{lemma} to \end{proof}
Let $z \colon [0, T] \to \mathbb{R}^4$ be a closed curve with
piecewise-constant velocity:
$\dot{z}(t) = w_k$ on $I_k = (\tau_{k-1}, \tau_k)$ for
$k = 1, \ldots, m$, where $0 = \tau_0 < \cdots < \tau_m = T$.
Then
\[
  \int_0^T \langle -J_0\, \dot{z}(t),\, z(t) \rangle \, dt
  = \sum_{1 \leq j < i \leq m}
    |I_j|\, |I_i|\, \omega_0(w_j, w_i).
\]
\end{lemma}

\begin{proof}
Since $z$ is closed, $\sum_{k=1}^m |I_k|\, w_k = 0$.
For $t \in I_i$:
\[
  z(t) = z(0) + \sum_{j=1}^{i-1} |I_j|\, w_j
    + (t - \tau_{i-1})\, w_i.
\]
Substitute into
$\langle -J_0 w_i, z(t) \rangle$ and integrate over $[0, T]$:
\begin{align*}
\int_0^T \!\langle{-}J_0 \dot{z}, z \rangle \, dt
  &= \sum_{i=1}^m \int_{I_i}
    \!\langle -J_0 w_i,\;
      z(0) + \textstyle\sum_{j < i} |I_j| w_j
      + (t{-}\tau_{i-1}) w_i
    \rangle\, dt \\
  &= \underbrace{
      \bigl\langle -J_0 \bigl(\textstyle\sum_i |I_i| w_i\bigr),\,
      z(0) \bigr\rangle
    }_{= \, 0 \;\text{(closure)}}
  + \sum_{i=1}^m \sum_{j=1}^{i-1}
      |I_i|\, |I_j|\, \omega_0(w_j, w_i) \\
  &\quad
  + \sum_{i=1}^m
      \underbrace{\omega_0(w_i, w_i)}_{= \, 0}\;
      \cdot \tfrac{|I_i|^2}{2}.
\end{align*}
The first term vanishes because
$\sum |I_k| w_k = 0$ (closure of~$z$).
The third term vanishes because
$\omega_0$ is antisymmetric.
\end{proof}


% ====================================================================
% HKO counterexample
% ====================================================================

\begin{theorem}[HKO counterexample to Viterbo's conjecture]%
\label{thm:hko-counterexample}
% [TODO: JÖRN - verify statement, add proof/reference]
% Source: known_polytopes.rs doc comment
% Reference: Haim-Kislev and Ostrover 2024
There exists a 10-facet polytope $K \subset \mathbb{R}^4$
(the HKO pentagon counterexample) with systolic ratio
$\operatorname{sys}(K) > 1$, disproving Viterbo's conjecture
in dimension~$4$.
\end{theorem}


\end{document}
